\documentclass{skripsi}
\usepackage{skripsi-telu-if}
\usepackage{booktabs}
\usepackage{lipsum}
\title[
en={Road Damage Prediction with Machine Learning}, 
id={Prediksi Kerusakan Jalan Raya dengan Machine Learning}
]{Prediksi Kerusakan Jalan Raya dengan Machine Learning}
\author[1234567890]{Asep Bensin}
\date{\today}
\firstadvisor[1234567890]{Dr.\ Ir.\ Ujang Sufahri, M.Eng.}
\secondadvisor[1234567890]{Prof.\ Dr.\ Bambang Supriyanto, M.Sc.}
\headdepartment[1234567890]{Dr.\ Eng.\ Siti Aminah, M.T.}
\addbibresource{final.bib}

\begin{document}
\frontmatter
\maketitle
\approvalpage
\originalitystatement
\onehalfspacing
\tableofcontents
\listoffigures
\listoftables
\begin{abstract}
\lipsum[1-2]
\end{abstract}
\mainmatter
\onehalfspacing
\chapter{Lorem}
\section{Consectetur}
\lipsum[1-3]
\subsection{Sit Amet}
\lipsum[4-6]
\subsubsection{Consectetur}
\lipsum[7-9]
\chapter{Ipsum}
\section{Amet}
\begin{figure}[h]
    \centering
    \includegraphics[width=0.5\textwidth]{example-image-b}
    \caption{Contoh Gambar}
    \label{fig:contoh-gambar}
\end{figure}
\lipsum[10-12]
\begin{table}[H]
    \centering
    \caption{Parameter-parameter dalam pencarian \textit{grid} tahap pertama (\textit{coarse grid-search}) untuk optimasi \textit{hyperparameter} model SVM.}
    \label{tab:grid_search_parameters}
    \begin{tabular}{lll}
        \toprule
        Parameter & Nilai yang Diuji & Jumlah Nilai \\
        \midrule
        % kernel
        kernel & RBF & 1 \\
        % regularization parameter
        $C$ & $\left\{ 2^{\,x} \,\middle|\, x \in \{-5, 0, 5, 10, 15\} \right\}$ & 5 \\
        $\gamma$ & $\left\{ 2^{\,x} \,\middle|\, x \in \{-15, -10, -5, 0, 5\} \right\}$ & 5 \\
        $n_{components}$ & $\{512, 256, 128, 64, 32, 16, 8, 4\}$ & 8 \\
        \midrule
        Total Kombinasi & & 200 \\
        \bottomrule
    \end{tabular}
\end{table}
\nocite{*}
\printbibliography[heading=bibintoc]

\end{document}